% Part: second-order-logic
% Chapter: syntax-and-semantics
% Section: introduction

\documentclass[../../../include/open-logic-section]{subfiles}

\begin{document}

\olfileid{sol}{syn}{int}

\olsection{परिचय}

प्रथम-क्रम तर्कशास्त्रात,
दिलेल्या भाषेतील
अतार्किक चिन्हे,
म्हणजे तिचे
!!{constant}s,
!!{function}s
आणि !!{predicate}s,
तार्किक चिन्हांबरोबर
जोडून आपण
प्रथम-क्रम
!!{structure}s
विषयी विधाने
मांडतो.
हे पूर्ति
या संकल्पनेने
केले जाते.
ती !!a{structure}~
$\Struct{M}$,
चर-मूल्यांकन~$s$
आणि !!a{formula}~$!A$
यांना जोडते:
$\Sat{M}{!A}[s]$
खरे असते, जर
आणि फक्त जर
त्यातील
!!{constant}s,
!!{function}s
आणि !!{predicate}s
यांचा अर्थ
$\Struct{M}$ नुसार
आणि त्यातील मुक्त
चलांचा अर्थ~$s$
नुसार घेतल्यावर
$!A$ जे व्यक्त
करते ते सत्य
असते.
!!{identity}~$\eq$
चा अर्थ
$\Sat{M}{!A}[s]$
च्या व्याख्येतच
दिलेला असतो;
$\lforall$ आणि
$\lexists$ यांचाही
अर्थ तसाच
दिलेला असतो.
यांपैकी पहिले
चिन्ह रचनेच्या
!!{domain}~$\Domain{M}$
वरील एकरूपता
संबंध म्हणूनच
नेहमी अर्थ लावले
जाते; संख्यकांचा
आवाका नेहमी
संपूर्ण !!{domain}
असतो.
पण महत्त्वाचे
म्हणजे संख्यकीकरण
फक्त !!{domain}
मधील घटकांवर
करता येते;
म्हणून संख्यकानंतर
फक्त वस्तू-!!{variable}s
येऊ शकतात.

द्वितीय-क्रम
तर्कशास्त्रात
भाषा आणि पूर्तीची
व्याख्या या दोन्ही
विस्तारून त्यांत
मुक्त व बद्ध
फलन-चल, विधेय-चल
आणि त्यांवरील
संख्यकीकरण
समाविष्ट केले
जाते.
वस्तू-चलांचा
!!{constant}s शी
जसा संबंध असतो,
तसाच या चलांचा
!!{function}s
आणि !!{predicate}s
शी असतो.
द्वितीय-क्रम
तर्कशास्त्रातील
पदे व !!{formula}s
घडवताना त्यांची
तशीच भूमिका
असते आणि
त्यांवरील
संख्यकीकरणही
तत्सम रीतीने
हाताळले जाते.
\emph{मानक}
चिन्हार्थमीमांसेत
द्वितीय-क्रम
संख्यकांचा
आवाका योग्य
प्रकारच्या सर्व
संभाव्य वस्तूंवर
असतो (फलन-चलांसाठी
$\Domain{M}$ वरून~$\Domain{M}$
मध्ये जाणारी
$n$-स्थानी फलने,
आणि विधेय-चलांसाठी
$n$-स्थानी संबंध).
उदाहरणार्थ,
$\lforall[\Obj{v_0}][(\Obj{P^1_0}(\Obj{v_0}) \lor \lnot
  \Obj{P^1_0}(\Obj{v_0}))]$
हे प्रथम-क्रम
आणि द्वितीय-क्रम
या दोन्ही
तर्कशास्त्रांतील
सूत्र आहे;
पण दुसऱ्यामध्ये
$\lforall[\Obj{V^1_0}][\lforall[\Obj{v_0}][(\Obj{V^1_0}(\Obj{v_0}) \lor
    \lnot \Obj{V^1_0}(\Obj{v_0}))]]$
आणि
$\lexists[\Obj{V^1_0}][\lforall[\Obj{v_0}][(\Obj{V^1_0}(\Obj{v_0}) \lor
    \lnot \Obj{V^1_0}(\Obj{v_0}))]]$
यांचाही विचार
करता येतो.
त्यांत मुक्त
चले नसल्याने
ती द्वितीय-क्रम
तर्कशास्त्रातील
!!{sentence}s
आहेत.
येथे $\Obj{V^1_0}$
हे द्वितीय-क्रम
$1$-स्थानी
विधेय-चल आहे.
$\Obj{V^1_0}$
चे अनुमत
अर्थनिर्धारण
$1$-स्थानी
!!{predicate}
असलेल्या
$\Obj{P^1_0}$
ला देता येणाऱ्या
अर्थनिर्धारणासारखेच
असते, म्हणजे
ते~$\Domain{M}$
चे उपसंच असतात.
त्यांवर संख्यकीकरण
केल्यास
$\lforall[\Obj{v_0}][(\Obj{V^1_0}(\Obj v_0) \lor \lnot
  \Obj{V^1_0}(v_0))]$
हे $\Obj{V^1_0}$
ला~$\Domain{M}$
चा उपसंच मूल्य
म्हणून देण्याच्या
प्रत्येक पद्धतीत
किंवा किमान
एका पद्धतीत
खरे आहे असे
म्हटले जाते.
प्रत्येक संचात
दिलेला घटक
असतो किंवा
नसतो; म्हणून
दोन्ही वाक्ये
कोणत्याही
!!{structure}
मध्ये सत्य असतात.
\end{document}
